\section{Radix-Packed Retrieval}
\label{sec:proposal}
The construction represents segments from different records as digits
inside a coefficient. It uses direct encryptions of public radix weights
and retains the linear reply interface of Section~\ref{sec:prelim}.
The encoding, selector, and extraction must share the same public layout.
\subsection{Record Encoding and Decoding}
\label{sec:encoding}
Fix an ordered record \(m_i=(b_{i,0},\ldots,b_{i,\ell-1})\in\{0,1\}^{\ell}\).
Set \(b_{i,v}=0\) for \(v\ge\ell\). For \(h\in[J]\), define
\begin{equation}
\label{eq:segments}
 x_{i,h}=\sum_{v=0}^{\rho_0-1}2^v b_{i,(h-1)\rho_0+v},\qquad
 0\le x_{i,h}<B=2^{\rho_0}.
\end{equation}
Set \(x_{i,h}=0\) for \(h>J\). Thus an incomplete final segment is padded
with zero high bits. Define \(\mathsf{Encode}(m_i)=(P_{i,1},\ldots,P_{i,L})\),
where
\begin{equation}
\label{eq:encode}
 P_{i,j}(X)=\sum_{u=0}^{n-1}x_{i,(j-1)n+u+1}X^u\in R_t,
 \qquad j\in[L].
\end{equation}
The indices are \(i\in[N]\), \(j\in[L]\), \(h\in[J]\), and
\(u\in\{0,\ldots,n-1\}\); their relation is \(h=(j-1)n+u+1\).
Each coefficient holds one scalar segment, and each polynomial holds up to
\(n\) segments. The final polynomial is padded with zero coefficients.

\(\mathsf{Decode}\) reads canonical coefficients in increasing \(j\), then
increasing \(u\), keeps the first \(J\) segments, and expands each into
\(\rho_0\) bits in increasing bit position. It concatenates these bits and
retains the first \(\ell\). For the image of \(\mathsf{Encode}\), uniqueness
of the binary expansion in \([0,B)\) recovers every bit in
\eqref{eq:segments}; the discarded bits and coefficients are precisely the
padding. Consequently
\[
 \mathsf{Decode}(\mathsf{Encode}(m_i))=m_i.
\]

\subsection{Packed Reply Structure}
\label{sec:overview}
For target ordinal \(\nu\in[\alpha]\), associate \(i_\nu\) with
\(k_\nu=B^{\nu-1}\). The intended plaintext of reply block \(j\) is
\begin{equation}\label{eq:packed-plaintext}
 PM_j=\sum_{\nu=1}^{\alpha}k_\nu P_{i_\nu,j}\in R_t.
\end{equation}
The client decrypts each block, extracts the canonical radix digits,
and decodes the resulting record polynomials (Fig.~\ref{fig:system-overview}).
\begin{figure}[t]
\centering
\resizebox{\linewidth}{!}{%
\begin{tikzpicture}[
 font=\small,>=Latex,
 box/.style={draw,rounded corners,align=center,text width=4.9cm,inner sep=5pt},
 flow/.style={-{Latex},thick}]
\node[box,text width=11.1cm,fill=black!4] (pp) at (0,0)
 {Public \(pp\): rings, distributions, \(K\), and layout\\
 \(J=\lceil\ell/\rho_0\rceil\) segments; \(L=\lceil J/n\rceil\) blocks};
\node[box] (client) at (-3,-2.2)
 {\textbf{Client: private \(sk=s\), \(I\), \(\mathsf{st}\)}\\
 \(s\leftarrow\mathsf{KeyGen}(pp)\)\\
 \(Q_i\leftarrow\mathsf{Enc}_s(v_{I,i})\)\\
 Fresh randomness for each \(i\in[N]\)};
\node[box] (server) at (3,-2.2)
 {\textbf{Server: \(pp\), encoded database}\\
 \(m_i\mapsto(P_{i,1},\ldots,P_{i,L})\)\\
 \(\mathsf{Resp}_j=\sum_i\Absorb(P_{i,j},Q_i)\)\\
 \(j\in[L]\); public linear evaluation};
\node[box,text width=11.1cm] (ext) at (0,-5.2)
 {\textbf{Client extraction}\\
 \(PM_j=\mathsf{Dec}_s(\mathsf{Resp}_j)\), \(y_{j,u}=[X^u]PM_j\)\\
 Scalar \(\mathsf{Unpack}(y_{j,u};B,\alpha)\), then \(\mathsf{Decode}\) for each target};
\draw[flow] (pp.south -| client.north) -- (client.north);
\draw[flow] (pp.south -| server.north) -- (server.north);
\draw[flow] (client.south) -- ++(0,-0.65) -| (server.south)
 node[pos=0.25,fill=white] {\(Q\): \(N\) ciphertexts};
\draw[flow] (server.east) -- (6.1,-2.2) -- (6.1,-4.1)
 -- (3,-4.1) -- (ext.north -| server.south);
\node[fill=white] at (4.5,-4.1) {\(\mathsf{Resp}\): \(L\) ciphertexts};
\end{tikzpicture}%
}
\caption{SKE query--reply--extraction flow. Secret-key generation and client
state are local to the client. Each reply encrypts one packed polynomial
\(PM_j\), containing up to \(n\) packed scalar coefficients; decryption
and extraction are correct under Theorem~\ref{thm:correctness}.}
\label{fig:system-overview}
\end{figure}
\subsection{Scalar Packing and Unpacking}
\label{sec:param}
\begin{lemma}[Radix packing/unpacking]
\label{lemma:packing-condition}
Let \(B=2^{\rho_0}\), \(\rho_0\ge1\), \(\alpha\ge1\), and
\(x_{\nu}\in\{0,\ldots,B-1\}\). Set
\begin{equation}
\label{eq:weights}
 K=(1,B,\ldots,B^{\alpha-1}),\qquad k_{\nu}=B^{\nu-1}.
\end{equation}
Then \(y=\sum_{\nu=1}^{\alpha}B^{\nu-1}x_{\nu}\) satisfies
\(0\le y<B^\alpha\), has a unique base-\(B\) representation, and
\(\mathsf{Unpack}(y;B,\alpha)\) returns its digits by
\begin{equation}
\label{eq:unpack}
 x_{\nu}=\left\lfloor\frac{y}{B^{\nu-1}}\right\rfloor\bmod B,
 \qquad \nu\in[\alpha].
\end{equation}
All such integers embed without reduction in \(\mathbb Z_t\) exactly when
\(t\ge B^\alpha\). For \(t=2^w\), this is equivalent to
\(w\ge\alpha\rho_0\).
\end{lemma}
\begin{proof}
The geometric sum gives
\[
 0\le y\le(B-1)\sum_{\nu=1}^{\alpha}B^{\nu-1}=B^\alpha-1.
\]
For each \(\nu\), the contribution from lower positions is less than
\(B^{\nu-1}\), while each higher position is a multiple of \(B^{\nu}\).
Integer division by \(B^{\nu-1}\), followed by reduction modulo \(B\),
therefore returns exactly \(x_{\nu}\). This deterministic inverse proves
uniqueness. Conversely, every integer in \([0,B^\alpha)\) has these
\(\alpha\) digits, so the domain has exactly \(B^\alpha\) values.
Its canonical embedding into \([0,t)\) is possible exactly when
\(t\ge B^\alpha\), including equality. When \(\alpha=1\),
\(K=(1)\) and the decoder returns \(x_1=y\).
\end{proof}

The weights need not be units modulo \(t\). Positional digit recovery
therefore needs no coprimality condition on individual weights, unlike
growing-weight remainder recovery~\cite{mulityquery}.
\subsection{Query Generation}
\label{sec:qgen}
Define the constant polynomial
\begin{equation}
\label{eq:selector}
 v_{I,i}=\begin{cases}B^{\nu-1},&i=i_{\nu}\text{ for some }\nu\in[\alpha],\\
                         0,&\text{otherwise}.\end{cases}
\end{equation}
It is a constant polynomial, not a polynomial with every coefficient equal
to \(k_{\nu}\). Algorithm~\ref{alg:qgen} calls encryption exactly once per
coordinate. It neither reuses one encryption of zero nor replaces direct
\(\mathsf{Enc}_s(k_{\nu})\) by \(k_{\nu}\mathsf{Enc}_s(1)\).
\begin{algorithm}[t]
\caption{Query Generation \(\mathsf{QueryGen}(pp,sk,I)\)}
\label{alg:qgen}
\begin{algorithmic}[1]
\Require Admissible \(pp\), client key \(sk=s\), legal ordered target tuple \(I\).
\Ensure Query \(Q\), private client state \(\mathsf{st}\).
\For{\(i=1\) to \(N\)}
 \State Determine \(v_{I,i}\) by \eqref{eq:selector}.
 \State \(Q_i\leftarrow\mathsf{Enc}_s(v_{I,i})\) using fresh independent \(a_i,e_i\).
\EndFor
\State \(\mathsf{st}\gets I\).
\State \Return \(((Q_1,\ldots,Q_N),\mathsf{st})\).
\end{algorithmic}
\end{algorithm}

\subsection{Reply Generation}
\label{sec:reply}
Algorithm~\ref{alg:replygen} uses the encoded polynomials. One ciphertext is
returned per polynomial block, regardless of whether its final coefficients
are padding. Only the client decrypts the response.
\begin{algorithm}[t]
\caption{Reply Generation \(\mathsf{ReplyGen}(pp,\{P_{i,j}\},Q)\)}
\label{alg:replygen}
\begin{algorithmic}[1]
\Require Public \(pp\), encoded database \(P_{i,j}\) for \(i\in[N],j\in[L]\), query \(Q\).
\Ensure Exactly \(L\) reply ciphertexts.
\For{\(j=1\) to \(L\)}
 \State \(\mathsf{Resp}_j\gets\sum_{i=1}^N\Absorb(P_{i,j},Q_i)\), using \(\mathsf{Add}\).
\EndFor
\State \Return \((\mathsf{Resp}_1,\ldots,\mathsf{Resp}_L)\).
\end{algorithmic}
\end{algorithm}

\subsection{Reply Extraction}
\label{sec:unpack}
Algorithm~\ref{alg:replyext} applies Lemma~\ref{lemma:packing-condition} to
each integer coefficient \(y_{j,u}\), then uses the record decoding already
defined in Section~\ref{sec:encoding}. Public \(pp\) fixes all lengths,
and derived radix weights. With \(\alpha=1,K=(1)\), the construction reduces to
the single-record reference flow in Section~\ref{sec:xpir-baseline}.
\begin{algorithm}[t]
\caption{Reply Extraction \(\mathsf{ReplyExt}(pp,sk,\mathsf{Resp},\mathsf{st})\)}
\label{alg:replyext}
\begin{algorithmic}[1]
\Require \(pp\), client key \(sk=s\), \(L\) reply ciphertexts, private \(\mathsf{st}=I\).
\Ensure Recovered tuple in the order of \(I\), under Theorem~\ref{thm:correctness}.
\For{\(j=1\) to \(L\)}
 \State \(PM_j\gets\mathsf{Dec}_s(\mathsf{Resp}_j)\).
 \For{\(u=0\) to \(n-1\)}
  \State \(y_{j,u}\gets[X^u]PM_j\), as an integer in \([0,t)\).
  \State \((\hat x_{1,h},\ldots,\hat x_{\alpha,h})\gets
   \mathsf{Unpack}(y_{j,u};B,\alpha)\), where \(h=(j-1)n+u+1\).
 \EndFor
\EndFor
\For{\(\nu=1\) to \(\alpha\)}
 \State Form \(\hat P_{\nu,j}=\sum_{u=0}^{n-1}\hat x_{\nu,(j-1)n+u+1}X^u\) for \(j\in[L]\).
 \State \(\hat m_{i_{\nu}}\gets\mathsf{Decode}(\hat P_{\nu,1},\ldots,\hat P_{\nu,L})\).
\EndFor
\State \Return \((\hat m_{i_1},\ldots,\hat m_{i_\alpha})\).
\end{algorithmic}
\end{algorithm}

\paragraph{Single-record reference.}
\label{sec:xpir-baseline}
With \(\alpha=1\) and \(K=(1)\), the query independently encrypts one at
the target and zero elsewhere. The reply contains \(L\) ciphertexts,
one for each record polynomial block. This is the flat XPIR-style
reference flow; native layouts outside this interface can have different
block counts.

\paragraph{Native representation.}
\label{sec:adapter}
An implementation must zero-extend each \(\rho_0\)-bit segment into its
declared \(w\)-bit field, pad to \(Ln\) fields, and import exactly the
polynomials in \eqref{eq:encode}. Extraction uses the same layout and
trims to \(\ell\) bits. Width changes require rebuilding preprocessing;
weights must be exactly representable by the encryption interface.
These representation conditions preserve the mathematics but do not
guarantee native support for every admissible parameter choice.
